\chapter{Discussion}
\label{ch:discussion}
\glsunset{ADNEX}
\glsunset{RMI}
\glsunset{IOTA}
\glsunset{AUC}
\glsunset{EIDM}
\glsunset{IPD}
\glsunset{SaMD}
\glsunset{MDR}
\glsunset{TRIPOD}
\glsunset{PROBAST}


This thesis set out to improve evidence-informed decision-making in the diagnosis of ovarian cancer by evaluating the \gls{ADNEX} model and advancing statistical methodology for \glspl{CPM}. The work integrates an applied arm --- focused on evaluating, comparing and updating \gls{ADNEX} in diverse clinical settings --- with a methodological arm that advances methods for model development, evaluation, and probability estimation. 

\section{Summary of findings}

Chapter~\ref{sec:chp_ADNEXSR} presented the first comprehensive systematic review and meta-analysis of \gls{ADNEX} external validation studies. Across 47 studies encompassing more than 17\,000 tumours, the model demonstrated robust discriminative performance, with a summary \gls{AUC} of 0.93 for distinguishing benign from malignant masses \cite{barrenadaADNEXRiskPrediction2024}. However, the review also exposed important weaknesses in the validation literature: 91\% of validations were at high risk of bias, primarily due to small sample sizes, inadequate handling of missing data, and incomplete performance assessment. Perhaps most critically, only 4 of 47 studies reported any form of calibration assessment, making it impossible to draw conclusions about the accuracy of \gls{ADNEX}'s predicted probabilities across settings.

Chapter~\ref{sec:chp_ADNEXvsRMI} extended this work with a head-to-head meta-analysis comparing \gls{ADNEX} to the \gls{RMI}, the most widely used alternative in clinical practice. The results were unequivocal: \gls{ADNEX} consistently outperformed \gls{RMI} in terms of discrimination (summary \gls{AUC} 0.92 vs 0.85 in operated patients) and clinical utility. The probability of \gls{ADNEX} being clinically useful in a hypothetical new centre was 96\%, compared to only 15\% for \gls{RMI} (cutoff of 200) at the 10\% risk threshold \cite{barrenadaHeadHeadComparison2025}. These findings provide strong evidence that clinical guidelines should recommend \gls{ADNEX} over \gls{RMI} for the preoperative characterisation of ovarian masses. The same conclusion was drawn from the recent Refining Ovarian Cancer Test Accuracy Scores (ROCkeTS) multicenter study \cite{sundarRiskpredictionModelsPostmenopausal2024,sundarDiagnosticTestsOvarian2026}.

A recurring theme throughout this thesis is that calibration --- the agreement between predicted probabilities and observed event rates --- remains underappreciated in the clinical prediction literature \cite{vancalsterCalibrationAchillesHeel2019}. The systematic reviews in Chapters~\ref{sec:chp_ADNEXSR} and~\ref{sec:chp_ADNEXvsRMI} showed that most validation studies of \gls{ADNEX} reported discrimination and classification but ignored calibration. That imbalance is methodologically problematic and clinically risky. A model can rank patients correctly and still provide systematically wrong absolute probabilities, which can be harmful when the output is used to decide whether a woman with an adnexal mass should be referred, conservatively managed, or undergo surgery \cite{vancalsterCalibrationRiskPrediction2015}. In the context of \gls{ADNEX} and prediction models as decision support tools, calibration and the probability estimate are not secondary outputs: they are the outputs that inform clinical decisions.

The same applied story motivates the most ambitious project, updating the ADNEX model to \emph{ADNEX2} (Chapter~\ref{sec:chp_ADNEX2}). The original \gls{ADNEX} development data were necessarily limited to women who underwent surgery, whereas in clinical practice the model is intended to be used to determine who to operate. We used Bayesian refitting to update \gls{ADNEX}, borrowing information from the original model and incorporating a more representative dataset that includes non-operated patients. The new \emph{ADNEX2} model is the first ovarian cancer diagnosis model developed using operated and conservatively managed patients. It has an AUC of 0.95 (0.93--0.96)  and it is better calibrated than the original ADNEX model. The new model is designed and intended to be used within the two-step strategy that starts with the \gls{BD}, followed by the model for indeterminate cases. 

Chapter~\ref{sec:chp_RFOverfitting} shifted attention from the applied model in ovarian cancer to the behaviour of algorithms used for probability estimation. Using both a case study in ovarian malignancy prediction and a simulation study, the chapter showed that random forests can achieve almost perfect training \gls{AUC}s, which is often a sign of overfitting, and at the same time competitive test performance. The key issue is not classical overfitting in the narrow sense, but the tendency of fully grown trees to create local probability peaks and unstable risk estimates. That distinction matters because prediction models are often tuned as if classification performance were the only objective. For risk prediction, however, the quality of the probabilities themselves is central, and the default strategy of growing trees as deep as possible can be counterproductive when the goal is calibration rather than only discrimination. We recommend not using the default hyperparameters for random forests when the goal is probability estimation, and instead tuning the tree depth and minimum node size to achieve better calibration \cite{barrenadaUnderstandingOverfittingRandom2024}.

Chapter~\ref{sec:chp_ClusteredCalibration} addressed a different but closely related issue: calibration should not be estimated assuming independence of observations  when the data come from multiple clusters. Much of the evidence used to evaluate \glspl{CPM} is clustered by hospital or study, yet traditional calibration plots suppress that structure. The chapter therefore proposed clustered group calibration (CG-C), two-stage meta-analysis calibration (2MA-C), and mixed model calibration (MIX-C) as solutions to present clustered calibration curves \cite{barrenadaClusteredFlexibleCalibration2026}. The main lesson is that accounting for clustering provides more accurate calibration plots. The recommended approaches, especially 2MA-C with splines and MIX-C, make that distinction visible by combining an overall calibration curve with prediction intervals where the calibration of a new centre should lie. That is a more realistic summary than a single pooled curve because it asks the clinically relevant question of whether a model remains well calibrated when it is moved to a new setting. The methods are implemented in the open source R package \texttt{ClusteredCalibration}.

Chapter~\ref{sec:chp_FundamentalProblem} took a step back from the technical evaluation of model performance to examine a more fundamental question: how much can we trust a risk estimate for an individual patient? The study illustrated that individual risk estimates are far more uncertain than commonly assumed. While estimation uncertainty (due to finite sample size) decreases with larger training samples, model uncertainty (due to arbitrary modelling choices) and applicability uncertainty (due to differences in measurement procedures and populations) remain substantial. Even with 10\,000 training observations, the 95\% range of predicted probabilities for the same patient averaged 0.39 on the probability scale \cite{barrenadaFundamentalProblemRisk2025}. These findings challenge the common practice of presenting a single risk estimate to patients and clinicians without acknowledging the underlying uncertainty.

Chapter~\ref{sec:chp_varselection} presents the first implementation of net benefit as a metric to optimise in variable selection. Across three case studies we present how this approach can lead to more parsimonious models with comparable or better clinical utility than traditional approaches. We then implement the algorithm as an open source R package (\texttt{NBvarsel}). This is a promising direction for future research, as it aligns the model development process with the ultimate goal of improving patient outcomes rather than just statistical performance.



\section{Embracing uncertainty in individualised decision-making}

\glspl{CPM} are developed and evaluated at the population level, but their purpose is to inform decisions for individual patients \cite{steyerbergClinicalPredictionModels2009}. The chapters in this thesis show that moving from population performance to individual decision-making adds successive layers of uncertainty. A patient and clinician may see one number, but that number is the product of a model that may have different performance in another centre, another patient subgroup, or another data collection context. Chapter~\ref{sec:chp_ClusteredCalibration} makes this visible at the centre level; Chapter~\ref{sec:chp_FundamentalProblem} makes it visible at the individual level. 

Chapters~\ref{sec:chp_ADNEXSR} and~\ref{sec:chp_ADNEXvsRMI} established that \gls{ADNEX} performs well on average across diverse populations. However, both chapters noted considerable heterogeneity in performance across centres and studies. This has direct implications for the concept of \gls{EIDM} introduced earlier in the thesis. A predicted probability should be treated as an estimate based on the used model, not as an exact measurement. Confidence intervals represent the model's uncertainty given the sample data, but they do not guarantee that a patient's true underlying risk lies within the interval, where \enquote{true underlying risk} is understood as the expected proportion of events in an infinite population of individuals sharing the exact same covariate pattern. In practice, that means risk estimates should be communicated together with an explanation of what is known, what is uncertain, and why the estimate may change if the patient were seen in another centre or assessed under slightly different assumptions \cite{pateUncertaintyUsingRisk2019, rileyUncertaintyRiskEstimates2025, aldermanTacklingAlgorithmicBias2025, ganapathiTacklingBiasAI2022}.

Chapter~\ref{sec:chp_FundamentalProblem} pushed this line of reasoning further, showing that even within a single centre, the predicted probability for an individual patient is subject to model and applicability uncertainty that cannot be reduced by increasing sample size alone \cite{pateUncertaintyUsingRisk2019, rileyUncertaintyRiskEstimates2025}. A risk estimate of 15\% should not be interpreted as a precise measurement but rather as a point within a range of plausible values, each of which might lead to a different clinical decision.

Our observations do not diminish the value of \glspl{CPM}. \gls{ADNEX} remains a useful tool for ovarian mass characterisation because it shows clinical utility at population level and can support referral and management decisions. The lesson is instead that its outputs should be used with caution. \glspl{CPM} should support shared decision-making rather than replace it, and that requires assessments embracing this uncertainty \cite{moonsCriticalAppraisalData2014, steyerbergClinicalPredictionModels2009}.

\section{Towards honest and safe medical devices}

The systematic reviews in Chapters~\ref{sec:chp_ADNEXSR} and~\ref{sec:chp_ADNEXvsRMI} consistently identified poor methodological quality in external validation studies. This is very concerning because this evidence is of utmost importance for safe implementation of \glspl{CPM} \cite{steyerbergPredictionModelsNeed2016}.  The most common issues --- assessed using the \gls{PROBAST} tool \cite{wolffPROBASTToolAssess2019} --- were inadequate sample sizes, inappropriate exclusion of patients with missing data, and incomplete reporting of model performance. Adherence to the \gls{TRIPOD} reporting guidelines \cite{collinsTRIPODAIStatementUpdated2024} was low, with key items such as sample size justification, missing data handling, and calibration assessment reported in fewer than 8\% of studies. These findings are consistent with previous reviews of clinical prediction model validation studies, which have repeatedly highlighted similar issues \cite{collinsExternalValidationMultivariable2014, bouwmeesterReportingMethodsClinical2012, andaurnavarroCompletenessReportingClinical2022,andaurnavarroSystematicReviewFinds2023, andaurnavarroSystematicReviewIdentifies2023, dhimanMethodologicalConductPrognostic2022, dhimanOverinterpretationFindingsMachine2023}.

This evidence is even more concerning in light of the regulatory developments discussed in the introduction. As \glspl{CPM} increasingly fall under medical device regulation --- including the EU \gls{MDR} and the \gls{SaMD} framework --- the evidence base for their validation must improve. Regulatory frameworks require demonstration of safety and effectiveness, which in the context of prediction models translates to robust evidence of discrimination, calibration, and clinical utility across diverse populations \cite{vancalsterPerformanceEvaluationPredictive2024, collinsTRIPODAIStatementUpdated2024}. The current state of the validation literature and regulatory approval falls short of what should be expected for tools that are intended to be used in patients. Recent studies have highlighted that many models that have been approved and are in use, have not undergone rigorous external validation, and their performance in real-world settings is often unknown or unclear \cite{mehtaEvaluatingTransparencyAI2025,vanleeuwenArtificialIntelligenceRadiology2021}.

Improving this situation requires combined efforts from researchers and regulators. Researchers should use better study designs, with adequate sample size \cite{rileyMinimumSampleSize2021}, proper handling of missing data and prospective multicenter data collection. The quality of reporting in published literature also needs improving, which can be done by adhering to international guidelines like \gls{TRIPOD}+AI and its extensions \cite{collinsTRIPODAIStatementUpdated2024,debrayTransparentReportingMultivariable2023, snellTransparentReportingMultivariable2023}, and a shift in culture towards recognising external validation as an ongoing process rather than a one-time exercise is needed \cite{vancalsterThereNoSuch2023}. From the regulators' side, there should be stricter requirements about the model's performance before approval, ensuring that at least key performance measures are presented, such as the calibration plot and a clinical utility metric (i.e. Net Benefit) \cite{vancalsterPerformanceEvaluationPredictive2024, ShowUsEvidence2026}. Both perspectives are synergistic: better research leads to better evidence for regulators, and stricter regulation incentivises better research. Ultimately, the goal is to ensure that \glspl{CPM} are safe, effective, and trustworthy tools for improving patient care.

\section{From publication to deployment}

The final aim of any prediction model should be to improve patient care. That requires moving beyond publication to deployment in clinical practice \cite{van2026enemies}. Developing and publishing a prediction model is a research endeavour, and we, researchers, are well trained to do that. Deploying a model in practice is more complicated. First it requires regulatory approval, which is a complex and often lengthy process that requires robust evidence of safety and effectiveness \cite{muehlematterApprovalArtificialIntelligence2021}. Second, it requires integration into clinical workflows, which can be challenging due to issues such as clinician acceptance, training, and changes in practice patterns \cite{kellyKeyChallengesDelivering2019}. Third, it requires ongoing monitoring of performance in real-world settings, which is essential to ensure that the model continues to perform well and does not cause unintended harms \cite{finlaysonClinicianDatasetShift2021}.  


This process requires a considerable amount of money and time, and the academic ecosystem is not prepared to support this process. Therefore, deployment is often left to industry, which will only invest in models with a clear path to commercialisation. This could generate unfair focus on models that are more likely to be profitable rather than those that are most needed or most effective. I believe that there should be incentives that support open source models that are developed in academia or non-profit organisations. Alternatively, close collaboration between academia and industry can ensure that the best models are deployed in practice, although this introduces conflicts of interest that can compromise trust and transparency \cite{pantanowitzRulesEngagementPromoting2022}. In these cases, boundaries between academia and industry need to be clearly defined beforehand.


For the developed model in this thesis, the \emph{ADNEX2} model, we closely collaborated with Gynaia Inc. (\url{https://www.gynaia.com/}), a KU Leuven spin-off that provides training and a compliant environment to implement the model. Currently \emph{ADNEX2} is undergoing medical device regulatory approval in the UK and the European Union, and will be deployed in the near future. 


\section{Future directions for trustworthy AI in medicine}

There is an exponential growth in the development of \glspl{CPM}, but the findings of this thesis suggest that there is still much work to be done to ensure that these models are trustworthy and useful in practice. A recent systematic review showed that the probability of a new clinical prediction model being externally validated in 5 years is only 13\% \cite{arshiNumberPublicationsNew2025}. Only one in eight models will be validated, which is a requirement for being trustworthy and applied in clinical practice. There is need for more and better validation studies, with a focus on calibration and clinical utility. The field of prediction models is moving fast, and there are many opportunities for future research to improve the development, evaluation, and implementation of \glspl{CPM}. Some promising directions include:

\subsection{Focus on Clinical Utility}

Clinical utility as measured by \gls{NB} or \gls{RU} is now mainly used to evaluate \glspl{CPM} performance \cite{vancalsterReportingInterpretingDecision2018}. In this thesis we proposed introducing \gls{NB} as a metric to optimise in variable selection and consequently only include variables in a model that have clinical utility (Chapter \ref{sec:chp_varselection}). Future research can expand this approach to other aspects of model development, such as hyperparameter tuning in machine learning algorithms. The metric could also be used directly in the loss function of the model, which would be a more direct way to optimise for clinical utility \cite{vickersHypothesisNetBenefit2025,chiLearningInterpretableRisk2026}. This is a promising direction for future research, as it aligns the model development process with the ultimate goal of improving patient outcomes rather than just statistical performance.

\subsection{Clustered calibration and uncertainty estimation}

The methods proposed in Chapter~\ref{sec:chp_ClusteredCalibration} for clustered calibration apply only to binary outcomes in diagnosis settings, but they could be extended to survival or competing risk models using similar rationale \cite{austinGraphicalCalibrationCurves2022,austinGraphicalCalibrationCurves2020,barrenadaClusteredFlexibleCalibration2026}. The current results highlighted a suboptimal estimation of the \glspl{PI}, where Bayesian approaches have been suggested to better account for the estimation uncertainty of the meta-analysis \cite{partlettRandomEffectsMetaanalysis2017}. Through the thesis we have observed that heterogeneity in performance across centres is common, and that it can have important implications for the interpretation of a model's performance and its generalisability; therefore it must be accounted for \cite{barrenadaADNEXRiskPrediction2024,barrenadaHeadHeadComparison2025}. Future research should continue to develop methods for clustered evaluation and meta-analysis with specific focus on accurate estimation of \glspl{PI}. Model evaluation should show not only the average performance but also the expected range of performance in a new centre, which is more relevant for clinical decision-making and regulatory approval \cite{vancalsterPerformanceEvaluationPredictive2024}.

% TODO update with reference to Ewout's update book and other references and set better guidance for future research in this area
% Another promising direction is to develop methods for local model update that can be applied when a model is moved to a new cluster. This could involve using the data from the new cluster to update the model parameters or to recalibrate the model, while borrowing information from the original development data \cite{jenkinsContinualUpdatingMonitoring2021, van2026enemies}. 
% These approaches could help to improve the performance of a model in a new setting and to better account for the uncertainty of predictions in that setting; however, the ultimate goal is to produce accurate individual risks. Therefore, future research should also focus on methods that provide \glspl{CI} that account for all sources of uncertainty in the prediction of an individual risk, including estimation uncertainty, model uncertainty and applicability uncertainty \cite{barrenadaFundamentalProblemRisk2025}. 

\subsection{Beyond logistic regression: new algorithms and data sources}
Most \glspl{CPM} are developed using logistic regression, which is a well-established and interpretable method that has been shown to even outperform more complex machine learning algorithms in many settings \cite{christodoulouSystematicReviewShows2019}. Later, gradient boosting methods, such as XGBoost, have become popular due to their strong performance for tabular data \cite{grinsztajnWhyDoTreeBased2022, chenXGBoostScalableTree2016}. The latest trend is the use of transformer-based tabular foundation models, which have shown impressive performance, outperforming any other model according to the well-known TabArena leaderboard for predictive machine learning \cite{quTabICLv2BetterFaster2026, hollmannAccuratePredictionsSmall2025, quTabICLTabularFoundation2025, ericksonTabArenaLivingBenchmark2025}. However, clinical data is no longer only tabular, but also includes images, volumes, audio, and multimodal data \cite{baltrusaitisMultimodalMachineLearning2017}. Future research should explore the use of these new algorithms and data sources for clinical prediction while taking into account that the model's quality is only going to be as good as the data used to train it \cite{Sanders2017GARBAGEIG}. 

\subsection{Bad big data or good small data?}

The availability of large open datasets has been a key driver of the vast number of published \glspl{CPM}, but it has also raised concerns about the quality and representativeness of these publications \cite{SPICK2026112203}. Open data is not a problem per se, but the provider of the data needs to present enough evidence to support the reliability of the data. A recent study highlighted this showing how two open datasets that were used for more than 125 publications --- even some patents --- had very poor quality and almost no information about data provenance \cite{gibsonEvidenceUnreliableData2026}. This is particularly concerning in the field of medicine, where the quality of data has direct implications for patient care. Ensuring that the data used for model development and validation is of high quality, representative of the target population, and collected in a transparent and ethical manner is essential for the development of trustworthy \glspl{CPM} \cite{collinsTRIPODAIStatementUpdated2024}. For \glspl{CPM}, the focus should be on collecting excellent quality data in multicenter studies under a clear shared protocol that will then allow building trustworthy models that can be safely implemented in practice \cite{collinsTRIPODAIStatementUpdated2024,moonsPROBASTAIUpdatedQuality2025,steyerberg2019clinical}. Ideally we would want good big data, but in the meantime, we should focus on using good small --- but sufficient --- data that will allow the development of useful models. In machine learning this concept is known as going from model-centric to data-centric AI, which is a promising direction for future research in \glspl{CPM} \cite{jakubikDataCentricArtificialIntelligence2024}.

\subsection{New players: policymakers and regulators as key stakeholders}

\gls{CPM} research has traditionally focused on two audiences: the technical audience including statisticians, engineers, data scientists and methodologists who develop and evaluate the models, and the clinical audience including clinicians who use the models in practice \cite{efthimiou_developing_2024}. However, from the moment the \glspl{CPM} are regulated as medical devices, policymakers and regulators are also key stakeholders in this field, as they are responsible for setting the standards for evidence generation and for approving the use of these models in practice \cite{fraserArtificialIntelligenceMedical2023}. In principle, \gls{MDR} is aligned with the goals of trustworthy AI, as it requires demonstration of safety, effectiveness, fairness and clinical utility across diverse populations \cite{vancalsterPerformanceEvaluationPredictive2024, collinsTRIPODAIStatementUpdated2024}. Future research should therefore also focus on engaging with policymakers and regulators to ensure that the state of the art recommendations in model development and validation are implemented in regulatory frameworks, and that regulators enforce these requirements in practice.  


\subsection{Is AI the future for ovarian cancer diagnosis?}

The work in this thesis has focused on the evaluation and update of the \gls{ADNEX} model for ovarian cancer diagnosis. This is a logistic multinomial model, which can be seen as very simple compared to the complex machine learning algorithms or the new transformer-based models that are being developed in the field of clinical prediction. However, the simplicity of the model is also one of its strengths, as it is interpretable, easy to implement, and has been shown to perform well across diverse settings \cite{barrenadaADNEXRiskPrediction2024}. Recently the IOTA7 multicenter study has finalized collection including more than 14\,000 patients, images and volumetric data, becoming the largest dataset ever collected for ovarian mass characterisation. This dataset will not only allow validation of the existing models for ovarian cancer diagnosis, including the \emph{ADNEX2} model developed in this thesis, but also allow exploration of new predictors, new modelling techniques and new data sources. Future research should also focus on the evaluation of the different models when used by examiners with different levels of expertise and to compare the performance of the model alone, the clinician assisted by the model and the clinician alone \cite{agarwalCombiningHumanExpertise}.


\section{Conclusion}
This thesis contributes to the development and evaluation of \glspl{CPM} for ovarian cancer diagnosis from both applied and methodological perspectives. The systematic evaluation of \gls{ADNEX} confirmed its value as a diagnostic tool and provided the evidence base for recommending it over alternatives such as the \gls{RMI}. At the same time, the methodological chapters showed why this should not be the end of the story: calibration is rarely assessed, probability estimates from AI prediction can be misleading if not validated properly, clustering changes how performance should be interpreted, and individual risk predictions carry more uncertainty than is often acknowledged.

These findings underscore the need for more research in the area of \glspl{CPM}. Although these models can be very valuable, they can also be harmful \cite{vancalsterCalibrationRiskPrediction2015}. Trustworthy AI requires honest model development and validation, ensuring that the results are transparent, reliable and generalisable. The \enquote{best model} in a paper should be the most useful one in practice too. Their responsible use requires rigorous validation, transparent communication of uncertainty, and ongoing monitoring --- principles that are increasingly formalised in both reporting guidelines and medical device regulation, although not properly enforced by some regulators \cite{mehtaEvaluatingTransparencyAI2025}. We must ensure that the rapid growth in the predictive AI field is matched by a commitment to rigorous evidence generation and responsible implementation, so that these tools can truly improve patient care without unintended harms \cite{ShowUsEvidence2026, goldenbergAIActuallyImproving2026}. 

\textbf{AI is shaping the field of medicine and ovarian cancer diagnosis, but we need to ensure that it is the right AI, and that it is used in the right way.}
